\section{正弦と余弦}\label{節:正弦余弦の交換}

前節の反転を区間 $\left[0,\frac{\pi}{2}\right]$ で使うと何が起こるだろうか.この区間での反転は,\,$x$ を $\frac{\pi}{2}-x$ に変える操作である.このとき
\[
\sin\left(\frac{\pi}{2}-x\right)=\cos x,\qquad
\cos\left(\frac{\pi}{2}-x\right)=\sin x
\]
なので,正弦と余弦が入れ替わる.この入れ替えが役立つのは,\,$q$ を $[0,1]$ 上の連続関数として
\[
I=\int_0^{\frac{\pi}{2}}
\frac{q(\sin x)}{q(\sin x)+q(\cos x)}\dd x
\]
の形の積分である.ただし,分母 $q(\sin x)+q(\cos x)$ は区間内で $0$ にならないとする.反転した積分ともとの積分を加えて被積分関数を定数 $1$ にできれば,積分は区間の長さだけで決まる.

反転すると,分子の $q(\sin x)$ は $q(\cos x)$ に変わり,分母は2つの項が入れ替わるだけで変わらない.したがって
\[
I=\int_0^{\frac{\pi}{2}}
\frac{q(\cos x)}{q(\sin x)+q(\cos x)}\dd x
\]
である.もとの式と加えれば,分子の和が分母そのものになる.
\[
\begin{aligned}
2I&=\int_0^{\frac{\pi}{2}}
\frac{q(\sin x)+q(\cos x)}{q(\sin x)+q(\cos x)}\dd x\\
 &=\int_0^{\frac{\pi}{2}}1\dd x\\
 &=\frac{\pi}{2}
\end{aligned}
\]
よって
\[
\int_0^{\frac{\pi}{2}}
\frac{q(\sin x)}{q(\sin x)+q(\cos x)}\dd x
 =\frac{\pi}{4}
\]
を得る.\,$q$ の具体的な形は一度も使っていない.重要なのは,反転で分子の2つの候補が交換されることなのだ.

\begin{bookpoint}
\textbf{POINT}\par
$x=\frac{\pi}{2}-t$ と置いて正弦と余弦を入れ替え,もとの積分と加えて分子の和を分母に一致させる.
\end{bookpoint}

\begin{example}{正弦と余弦の交換}
次の定積分を求めよ.
\[
I=\int_0^{\frac{\pi}{2}}
\frac{\sin^4x}{\sin^4x+\cos^4x}\dd x.
\]
\end{example}

\paragraphbf{方針}
本節の式で $q(s)=s^4$ とした場合である.反転で分子の $\sin^4x$ が $\cos^4x$ に変わり,2つの分子の和がちょうど分母 $\sin^4x+\cos^4x$ になる.

\paragraphbf{解答}
$I=\frac{\pi}{4}$

\paragraphbf{解法}
$\sin x$ と $\cos x$ は同時に $0$ にならないので,分母は区間全体で正である.\,$x=\frac{\pi}{2}-t$ と置くと,分子は $\cos^4t$ に変わり,分母は変わらない.積分変数の文字を $x$ に戻せば
\[
I=\int_0^{\frac{\pi}{2}}
\frac{\cos^4x}{\sin^4x+\cos^4x}\dd x
\]
である.2つを加えれば,分子の和が分母そのものになる.
\[
\begin{aligned}
2I&=\int_0^{\frac{\pi}{2}}1\dd x\\
 &=\frac{\pi}{2}
\end{aligned}
\]
よって $I=\frac{\pi}{4}$ である.これは上の一般形で $q(t)=t^4$ とした場合である.\,$\sin^4x$ の次数を下げたり,分母を展開したりする必要はない.

なお,一般形は $q(t)=t^4$ に限らない.たとえば $q(t)=t^2+1$ でも,分母は $\sin^2x+\cos^2x+2=3$ で $0$ にならず,同じ値 $\frac{\pi}{4}$ になる.

\paragraphbf{図形的解釈}
図\ref{図:相補重み}は,例題の被積分関数を $w(x)$ としたグラフである.左右を反転したグラフ $w(\frac{\pi}{2}-x)$ と,同じ $x$ における高さを足すと $1$ になる.2つのグラフの下の面積は等しく,その和は縦 $1$,横 $\frac{\pi}{2}$ の長方形の面積である.したがって $w(x)$ のグラフは,この長方形の面積を2等分しているのである.

\begin{figure}[htbp]
\centering
\begin{tikzpicture}[x=3.5cm,y=2.3cm,>=Stealth]
\fill[blue!10](0,0)--plot[domain=0:1.570796,samples=100](\x,{pow(sin(\x r),4)/(pow(sin(\x r),4)+pow(cos(\x r),4))})--(1.570796,0)--cycle;
\draw[->](-0.08,0)--(1.73,0)node[right]{$x$};
\draw[->](0,-0.08)--(0,1.2)node[above]{$y$};
\draw[dashed](0,1)--(1.570796,1)--(1.570796,0);
\draw[dashed](0.785398,0)--(0.785398,1);
\draw[blue!75!black,thick]plot[domain=0:1.570796,samples=100](\x,{pow(sin(\x r),4)/(pow(sin(\x r),4)+pow(cos(\x r),4))});
\draw[orange!80!black,thick,dashed]plot[domain=0:1.570796,samples=100](\x,{pow(cos(\x r),4)/(pow(sin(\x r),4)+pow(cos(\x r),4))});
\node[left]at(0,1){$1$};\node[below]at(0,0){$0$};
\node[below]at(0.785398,0){$\frac{\pi}{4}$};
\node[below]at(1.570796,0){$\frac{\pi}{2}$};
\node[blue!75!black]at(1.15,0.65){$w(x)$};
\node[orange!80!black]at(0.48,1.26){$w(\frac{\pi}{2}-x)$};
\end{tikzpicture}
\caption{反転した2つの重みは,同じ場所で足すと $1$ になる.}\label{図:相補重み}
\end{figure}
